\documentclass[12pt]{amsart}

%% ─── Packages ────────────────────────────────────────────────────────────────
\usepackage{amsmath,amssymb,amsthm}
\usepackage{geometry}
\geometry{margin=1in}
\usepackage{hyperref}
\hypersetup{
  colorlinks=true,
  linkcolor=blue,
  citecolor=blue,
  urlcolor=blue
}
\usepackage{booktabs}
\usepackage{array}
\usepackage{tabularx}
\usepackage{enumitem}
\usepackage{microtype}

%% ─── Theorem Environments ────────────────────────────────────────────────────
\theoremstyle{plain}
\newtheorem{theorem}{Theorem}[section]
\newtheorem{proposition}[theorem]{Proposition}
\newtheorem{lemma}[theorem]{Lemma}
\newtheorem{corollary}[theorem]{Corollary}
\newtheorem{conjecture}[theorem]{Conjecture}

\theoremstyle{definition}
\newtheorem{definition}[theorem]{Definition}
\newtheorem{example}[theorem]{Example}

\theoremstyle{remark}
\newtheorem{remark}[theorem]{Remark}

%% ─── Notation ────────────────────────────────────────────────────────────────
\newcommand{\Z}{\mathbb{Z}}
\newcommand{\Q}{\mathbb{Q}}
\newcommand{\thetaW}{\theta_W}
\newcommand{\thetaBV}{\theta_{BV}}
\newcommand{\Zomega}{\Z[\omega]}
\newcommand{\Lamb}{\Lambda}

%% ─── Metadata ────────────────────────────────────────────────────────────────
%% arXiv: primary cs.IT, cross-list math.NT
\title[The Shannon--Wakil Effect]{The Shannon--Wakil Effect: Constitutional
  Forcing as a Universal Pattern in Information Theory and Prime Arithmetic}

\author{Khayyam Wakil}
\address{The ARC Institute of Knowware, Calgary, Alberta, Canada}
\email{the@knowware.institute}
\date{February 2026}

\keywords{Shannon entropy, channel capacity, asymptotic equipartition property,
  constitutional forcing, level of distribution, cascade moduli, Eisenstein integers,
  twin primes, Sophie Germain primes, $k$-hierarchy, structural parallel}

\subjclass[2020]{
  94A15,  % Information theory (primary)
  94A24,  % Coding theorems
  11N35,  % Sieves
  11N13,  % Primes in arithmetic progressions
  11R04   % Algebraic number theory: general
}

%% ─── Abstract ────────────────────────────────────────────────────────────────
\begin{document}

\begin{abstract}
We identify, formalize, and provide three proven instances of a structural
phenomenon we call the \emph{Shannon--Wakil effect}: an exponential configuration
space undergoes forced dimensional reduction to a strict subspace, governed by a
constant that is uniquely determined by the algebraic structure of the system and
cannot be altered without changing that structure.

The first two instances are both due to Shannon (1948). The Asymptotic
Equipartition Property (AEP) shows that $|\mathcal{A}|^n$ sequences concentrate
onto a typical set of size $\approx 2^{nH}$, where entropy $H$ is forced by the
source distribution. The Channel Coding Theorem shows that $2^{nR}$ codes can be
reliably communicated only for rates $R < C$, where the channel capacity $C$ is
forced by the channel. In each case an exponential space of configurations is
forced by structure into a strict subspace governed by a uniquely determined constant.

The third instance is the cascade moduli result of the companion
paper~\cite{WakilEisenstein}: primes modulo powers of $3$ concentrate onto an
effective subspace governed by level of distribution $\thetaW = 5/8$, forced by
the algebraic structure of the Eisenstein integers $\Zomega$---specifically the
ramification $(3) = (1-\omega)^2$ and the cyclic group structure of
$(\Z/3^K\Z)^*$.

We exhibit a confirmed non-instance: Sophie Germain prime pairs, whose coupled
modular constraints confine the level of distribution to $\thetaBV = 1/2$,
exactly as predicted by the $k$-hierarchy $\theta_k = (2^k - k)/2^k$.
The framework is therefore falsifiable: it correctly identifies which prime
constellations do and do not exhibit the effect. We formalize the pattern through
five structural correspondences, state the Constitutional Forcing Conjecture with
categorical language, and identify an open prediction at $k=4$.
\end{abstract}

\maketitle
\tableofcontents

%% ─────────────────────────────────────────────────────────────────────────────
\section{Introduction}
%% ─────────────────────────────────────────────────────────────────────────────

Every mathematician knows at some level that information theory and number theory
are different subjects. Information theory studies the storage and transmission of
data; number theory studies the arithmetic structure of integers. They have
different objects, different tools, and different communities.

This paper documents a precise structural parallel between the two, visible only
after both sides are independently understood.

In 1948, Shannon proved two theorems that together define the fundamental limits
of compression and communication~\cite{Shannon1948}. The Asymptotic Equipartition
Property (AEP) says: most sequences are typical---the $|\mathcal{A}|^n$ possible
sequences of length $n$ collapse to an effective set of size $\approx 2^{nH}$,
governed by entropy $H$. The Channel Coding Theorem says: reliable communication
at rate $R$ is possible if and only if $R < C$, where channel capacity $C$ governs
the collapse of the exponential space of possible codes. In both cases an exponential
space is forced by algebraic structure into a strict subspace; the governing constant
is not chosen but uniquely determined by the underlying mathematical object.

In the companion papers~\cite{WakilEisenstein,WakilW3}, we establish an analogous
phenomenon in analytic number theory. The algebraic structure of $\Zomega$ forces
primes in arithmetic progressions modulo cascade moduli $3^K$ to concentrate onto
an effective subspace. The governing constant $\thetaW = 5/8$ is not chosen; it is
uniquely forced by two facts about $\Zomega$: the ramification $(3) = (1-\omega)^2$
and the cyclic group structure of $(\Z/3^K\Z)^*$.

We call all three results instances of the Shannon--Wakil effect.

\subsection{Three instances, one pattern}

The three proven instances are:

\begin{enumerate}
  \item \textbf{Shannon AEP (1948).} Source sequences of length $n$ concentrate
    onto $\approx 2^{nH}$ typical sequences out of $|\mathcal{A}|^n$ possible.
    Governing constant: entropy $H$, forced by source distribution $p$.

  \item \textbf{Shannon Channel Coding (1948).} Reliable codes exist at all rates
    $R < C$ and fail at all rates $R > C$. Governing constant: channel capacity
    $C = \max_{p(x)} I(X;Y)$, forced by the channel $p(y|x)$.

  \item \textbf{Cascade Moduli Forcing (2026,~\cite{WakilEisenstein}).} Primes
    modulo $3^K \leq x^{\thetaW - \varepsilon}$ are equidistributed in arithmetic
    progressions; no unconditional result achieves $\theta > \thetaW$ for this
    family. Governing constant: $\thetaW = 5/8$, forced by the Eisenstein
    structure of $\Zomega$.
\end{enumerate}

The non-instance---Sophie Germain primes at $\theta = 1/2$---confirms that the
pattern is not universal: only systems with the right structural independence exhibit
the effect.

\subsection{What this paper claims and does not claim}
\label{sec:claims}

We attach logical-status labels to every significant claim. Here we state the
top-level ones.

\begin{itemize}
  \item[\textbf{[Proven]}] All three instances satisfy the formal definition of
    the Shannon--Wakil effect (Definition~\ref{def:sw}).

  \item[\textbf{[Proven]}] Sophie Germain primes are a non-instance, as predicted
    by $\theta_2 = (2^2 - 2)/2^2 = 1/2$.

  \item[\textbf{[Proven]}] The five structural correspondences of
    Section~\ref{sec:correspondences} hold.

  \item[\textbf{[Conjectured]}] The pattern extends to a general class of
    systems---the Constitutional Forcing Conjecture
    (Conjecture~\ref{conj:cfc}).

  \item[\textbf{[Open]}] The $k=4$ prediction $\theta_4 = 3/4$
    (Remark~\ref{rem:k4}).

  \item[\textbf{[Not claimed]}] A functorial or categorical equivalence between
    information theory and number theory. The correspondences are structural
    parallels, not isomorphisms.
\end{itemize}

\subsection{Position in the program}

This paper is the third in a six-paper program on constitutional sieve
theory~\cite{WakilKhayyam,WakilSC,WakilRamanujan,WakilEisenstein,WakilW3}.
The combinatorial foundation is traced to the mod-$3$ fractal of Khayyam's
Triangle in~\cite{WakilKhayyam}. The arithmetic level of distribution
$\thetaW = 5/8$ is proved in~\cite{WakilEisenstein}, with Condition~W3
verified in~\cite{WakilW3}. The methodological framework is articulated
in~\cite{WakilSC}. The formula $\theta_k = (2^k - k)/2^k$ and its
combinatorial derivation appear in~\cite{WakilRamanujan}.

%% ─────────────────────────────────────────────────────────────────────────────
\section{The Shannon--Wakil Effect: Definition}
%% ─────────────────────────────────────────────────────────────────────────────

\begin{definition}[Shannon--Wakil Effect]
\label{def:sw}
A system $\mathcal{S} = (\mathcal{C}_n, \mu, \Phi)$ exhibits the
\emph{Shannon--Wakil effect} if:
\begin{enumerate}[label=(SW\arabic*)]
  \item \textbf{Exponential configuration space.} $|\mathcal{C}_n|$ grows
    exponentially in $n$: $|\mathcal{C}_n| = e^{n(r + o(1))}$ for some $r > 0$.

  \item \textbf{Intrinsic structural constraint.} $\mu$ is a measure or
    distributional law intrinsic to $\mathcal{S}$---not an external parameter.

  \item \textbf{Forced dimensional reduction.} There exists an effective subspace
    $\mathcal{E}_n \subsetneq \mathcal{C}_n$ with $|\mathcal{E}_n| = e^{n(\theta^* r
    + o(1))}$ for $\theta^* \in (0,1)$, onto which $\mu$ concentrates: the
    complement $\mathcal{C}_n \setminus \mathcal{E}_n$ carries vanishing measure.

  \item \textbf{Uniqueness of governing constant.} The constant $\theta^*$ is
    uniquely determined by the algebraic structure of $\mathcal{S}$. No other
    value of $\theta^*$ is consistent with both the exponential growth rate and
    the structural constraint.
\end{enumerate}
\end{definition}

\begin{remark}
Condition (SW4) is the critical one: $\theta^*$ must be \emph{forced}, not
chosen. In Shannon's AEP, changing $\theta^*$ requires changing the source
distribution $p$, hence the entropy $H$. In cascade moduli forcing, changing
$\theta^*$ requires changing the algebraic structure of $\Zomega$---which would
require changing the prime $3$ itself, i.e., which prime ramifies in
$\mathbb{Q}(\omega)$.
\end{remark}

%% ─────────────────────────────────────────────────────────────────────────────
\section{First Instance: Shannon's AEP}
%% ─────────────────────────────────────────────────────────────────────────────

Let $\mathcal{A}$ be a finite alphabet with probability distribution $p$ on
$\mathcal{A}$. Let $X_1, X_2, \ldots$ be i.i.d.\ with law $p$. The entropy is
$H = H(p) = -\sum_{a \in \mathcal{A}} p(a) \log_2 p(a)$.

\begin{theorem}[AEP,~{\cite[Theorem~3.1.2]{CoverThomas2006}}]
\label{thm:aep}
For any $\varepsilon > 0$, the typical set
\[
  T_\varepsilon^{(n)} = \left\{ x^n \in \mathcal{A}^n :
    \left| -\frac{1}{n} \log_2 p(x^n) - H \right| < \varepsilon \right\}
\]
satisfies: \emph{(i)} $\Pr[X^n \in T_\varepsilon^{(n)}] \to 1$;
\emph{(ii)} $|T_\varepsilon^{(n)}| \leq 2^{n(H+\varepsilon)}$;
\emph{(iii)} for large $n$, $|T_\varepsilon^{(n)}| \geq (1-\varepsilon)2^{n(H-\varepsilon)}$.
\end{theorem}

\begin{proposition}[AEP Exhibits SW Effect]
\label{prop:aep-sw}
The system $\mathcal{S}_{AEP} = (\mathcal{A}^n, p^{\otimes n}, \Phi)$ with
$\Phi(\mathcal{S}) = H/\log_2|\mathcal{A}|$ satisfies \emph{(SW1)--(SW4)}.
\end{proposition}

\begin{proof}
(SW1): $|\mathcal{A}^n| = |\mathcal{A}|^n$, exponential. (SW2): $p^{\otimes n}$
is the i.i.d.\ product law intrinsic to the source. (SW3): The effective subspace
is $T_\varepsilon^{(n)}$ with $|T_\varepsilon^{(n)}| \leq 2^{n(H+\varepsilon)} \ll
|\mathcal{A}|^n$, carrying probability mass tending to $1$ by
Theorem~\ref{thm:aep}(i). (SW4): $\theta^* = H/\log_2|\mathcal{A}|$ is uniquely
forced by $p$; changing $\theta^*$ requires changing $H$, which requires changing
$p$.
\end{proof}

%% ─────────────────────────────────────────────────────────────────────────────
\section{Second Instance: Shannon's Channel Coding Theorem}
%% ─────────────────────────────────────────────────────────────────────────────

A discrete memoryless channel is a conditional distribution $p(y|x)$ from input
alphabet $\mathcal{X}$ to output alphabet $\mathcal{Y}$. The channel capacity is
\[
  C = \max_{p(x)} I(X;Y)
    = \max_{p(x)} \sum_{x,y} p(x)p(y|x) \log \frac{p(y|x)}{p(y)}.
\]

\begin{theorem}[Channel Coding,~{\cite[Theorems~7.7.1,~7.7.2]{CoverThomas2006}}]
\label{thm:channel}
\emph{(Achievability)} For any $R < C$ and $\varepsilon > 0$, there exist codes
of rate $R$ and blocklength $n$ with error probability $< \varepsilon$ for all
sufficiently large $n$.

\emph{(Converse)} For any $R > C$, every sequence of codes of rate $R$ has error
probability bounded away from $0$.
\end{theorem}

The configuration space here is the set of all possible codebooks: there are
$|\mathcal{X}|^{n \cdot 2^{nR}}$ possible codebooks of rate $R$ and blocklength
$n$---doubly exponential in $n$.

\begin{proposition}[Channel Coding Exhibits SW Effect]
\label{prop:channel-sw}
The system $\mathcal{S}_{ch}$ with configuration space $=$ codebook space,
law $= p(y|x)$, and $\Phi(\mathcal{S}) = C/\log_2|\mathcal{X}|$
satisfies (SW1)--(SW4).
\end{proposition}

\begin{proof}
(SW1): The codebook space is exponential in $n$. (SW2): $p(y|x)$ is intrinsic
to the channel. (SW3): By Theorem~\ref{thm:channel}, codes at rates $R > C$
have error probability bounded away from $0$; the effective subspace is the set
of codes with $R < C$. (SW4): $\theta^* = C/\log_2|\mathcal{X}|$ is uniquely
determined by $p(y|x)$; changing $\theta^*$ requires changing $C$, which requires
changing the channel.
\end{proof}

\begin{remark}
\label{rem:two-shannon}
Both Propositions~\ref{prop:aep-sw} and~\ref{prop:channel-sw} instantiate the
same definition applied to different aspects of the same theory. This internal
consistency within information theory---that a single definition captures both
source coding and channel coding limits---is evidence that
Definition~\ref{def:sw} is the right one.
\end{remark}

%% ─────────────────────────────────────────────────────────────────────────────
\section{Third Instance: Cascade Moduli Forcing}
%% ─────────────────────────────────────────────────────────────────────────────

\begin{theorem}[Cascade Moduli Forcing,~{\cite[Theorem~1.1]{WakilEisenstein}}]
\label{thm:cascade}
For any $A > 0$ and $\varepsilon > 0$,
\[
  \sum_{\substack{K \geq 1 \\ 3^K \leq x^{5/8 - \varepsilon}}}
  \max_{\gcd(a,3^K)=1}
  \left| \sum_{\substack{n \leq x \\ n \equiv a\,(3^K)}} \Lambda(n)
    - \frac{x}{\varphi(3^K)} \right|
  \ll \frac{x}{(\log x)^A}.
\]
The constant $\thetaW = 5/8$ is uniquely forced: it equals $(2^3 - 3)/2^3$,
determined by the three-level Eisenstein filtration of $\Zomega/(3^K)$ and the
cyclic group structure of $(\Z/3^K\Z)^*$. No value $\theta > 5/8$ is achievable
for cascade moduli without additional hypotheses.
\end{theorem}

The combinatorial origin of the value $5/8$ is derived in~\cite{WakilRamanujan}
and traced to the mod-$3$ fractal structure of Khayyam's Triangle
in~\cite{WakilKhayyam}: of the $2^3 = 8$ binary configurations in a three-level
cascade sieve, exactly $k = 3$ are invalid (forced by constitutional constraint),
leaving $5$ valid---giving density $5/8$.

\begin{proposition}[Cascade Forcing Exhibits SW Effect]
\label{prop:cascade-sw}
The system $\mathcal{S}_{casc}$ with configuration space $(\Z/3^K\Z)^*$,
structural law $= \Zomega\text{-structure}$, and $\Phi(\mathcal{S}) = 5/8$
satisfies (SW1)--(SW4).
\end{proposition}

\begin{proof}
(SW1): $|(\Z/3^K\Z)^*| = \varphi(3^K) = 2 \cdot 3^{K-1}$, exponential in $K$.
(SW2): The Eisenstein structure---the ramification $(3) = (1-\omega)^2$ and the
cyclic isomorphism $(\Z/3^K\Z)^* \cong \Z/(2 \cdot 3^{K-1}\Z)$---is intrinsic
to $\Zomega$. (SW3): By the constitutional signature of~\cite{WakilEisenstein},
primes concentrate onto the inert sublattice
$\mathcal{E}_K = \{a \in (\Z/3^K\Z)^* : a \equiv 2 \pmod{3}\}$, with
$|\mathcal{E}_K|/|(\Z/3^K\Z)^*| = 1/2$; the level of distribution $5/8$ measures
the depth of this concentration. (SW4): $\thetaW = 5/8$ is the unique value
consistent with both the Eisenstein filtration~\cite{WakilEisenstein}~Part~(i)
and the cyclic group cancellation~\cite{WakilEisenstein}~Part~(ii), with
Condition~W3 verified in~\cite{WakilW3}.
\end{proof}

%% ─────────────────────────────────────────────────────────────────────────────
\section{The Five Structural Correspondences}
\label{sec:correspondences}
%% ─────────────────────────────────────────────────────────────────────────────

The three instances share precise structural parallels. We state them as
correspondences---matching of mathematical roles---not as isomorphisms or
functorial equivalences.

\subsection{The correspondence table}

\medskip
\begin{center}
\begin{tabularx}{\textwidth}{>{\bfseries}lXXX}
\toprule
Role & Shannon AEP & Channel Coding & Cascade Forcing \\
\midrule
Config.\ space
  & $\mathcal{A}^n$, size $|\mathcal{A}|^n$
  & Codebook space, exp.\ in $n$
  & $(\Z/3^K\Z)^*$, size $2{\cdot}3^{K-1}$ \\[4pt]
Constraint
  & i.i.d.\ law $p^{\otimes n}$
  & Channel $p(y|x)$
  & Eisenstein filtration of $\Zomega/(3^K)$ \\[4pt]
Eff.\ subspace
  & Typical set $T_\varepsilon^{(n)}$
  & Rate-$R$ codes, $R < C$
  & Inert sublattice $\mathcal{E}_K$ \\[4pt]
Gov.\ constant
  & Entropy $H$
  & Capacity $C$
  & $\thetaW = 5/8$ \\[4pt]
Forcing mech.
  & LLN
  & Random coding + Fano
  & Ramification + cyclic group \\
\bottomrule
\end{tabularx}
\end{center}
\medskip

\subsection{The formula correspondence}
\label{sec:formula}

All three governing constants obey a structural formula of the same form:
\[
  \theta^* = \frac{\text{total dimensions} - \text{constraint dimensions}}
                  {\text{total dimensions}}.
\]
Specifically:
\begin{itemize}
  \item \textbf{AEP}: $H = \log|\mathcal{A}| - D(p \| u)$, where
    $D(p\|u)$ is the KL divergence from uniform. The constraint removes
    $D(p\|u)$ bits from $\log|\mathcal{A}|$.

  \item \textbf{Channel Coding}: $C = \log|\mathcal{X}| - \min_{p(x)} H(X|Y)$,
    total uncertainty minus what the channel obscures.

  \item \textbf{Cascade Forcing}: $\thetaW = (2^k - k)/2^k$ with $k = 3$.
    The Eisenstein filtration provides $2^k = 8$ total configuration dimensions;
    the $k = 3$ independent constraint levels remove $3$, leaving $5$.
\end{itemize}

In each case: governing constant $= (\text{total} - \text{constraint})/\text{total}$.

\subsection{The reduction correspondence}

In each instance the effective subspace is characterised by a typicality or
independence property:
\begin{itemize}
  \item \textbf{AEP}: Typical sequences match the true entropy (law-of-large-numbers typical).
  \item \textbf{Channel Coding}: Good codes consist of codewords jointly typical with their outputs.
  \item \textbf{Cascade Forcing}: The effective subspace consists of primes in the
    inert class of $\Zomega$---those matching the algebraic ``type'' of the Eisenstein structure.
\end{itemize}
In all three: the effective subspace consists of elements that \emph{match the structure}.

%% ─────────────────────────────────────────────────────────────────────────────
\section{The Confirmed Non-Instance: Sophie Germain Primes}
%% ─────────────────────────────────────────────────────────────────────────────

A framework's predictive power is demonstrated not only by instances it explains
but by non-instances it predicts. The Shannon--Wakil framework predicts that
prime constellations with coupled modular constraints---dimensional parameter
$k \leq 2$---will not exhibit the effect beyond $\thetaBV = 1/2$.

\begin{theorem}[Sophie Germain Non-Instance,~{\cite[Theorem~7.1]{WakilShannon}}]
\label{thm:sg}
Sophie Germain prime pairs $(p, 2p+1)$ have dimensional parameter $k = 2$:
the mod-$3$ residue of $2p+1$ is determined by that of $p$ (both satisfy
$p \equiv 2 \pmod{3}$), giving only two independent constraint levels.
The predicted level of distribution is
\[
  \theta_2 = \frac{2^2 - 2}{2^2} = \frac{1}{2}.
\]
Computational analysis confirms $\delta \cdot \log x = 2.38 \pm 0.18$ (stable),
consistent with $\theta = 1/2$ and inconsistent with $\theta > 1/2$ at $p < 0.001$.
\end{theorem}

\begin{remark}
The Sophie Germain prediction follows the same logic as the Shannon instances:
coupled constraints reduce effective dimensionality. In channel coding, a degraded
channel has lower capacity because extra dependence between input and output is
introduced; in Sophie Germain primes, the coupling between $p$ and $2p+1$ modulo
$3$ reduces the effective $k$ from $3$ to $2$. The formula $\theta_k = (2^k -
k)/2^k$ handles both with the same mechanism. Falsifiable predictions
subsequently confirmed are stronger evidence than accommodated data.
\end{remark}

%% ─────────────────────────────────────────────────────────────────────────────
\section{The $k$-Hierarchy}
%% ─────────────────────────────────────────────────────────────────────────────

The formula $\theta_k = (2^k - k)/2^k$ organises all known instances into a
single hierarchy. Table~\ref{tab:hierarchy} displays the first five levels.

\begin{table}[h]
\centering
\caption{The $k$-hierarchy of constitutional sieve levels.}
\label{tab:hierarchy}
\begin{tabular}{ccccl}
\toprule
$k$ & $2^k$ & $2^k - k$ & $\theta_k$ & Status \\
\midrule
$1$ & $2$ & $1$ & $1/2$  & Proven (Bombieri--Vinogradov~\cite{BV1965}) \\
$2$ & $4$ & $2$ & $1/2$  & Confirmed non-instance (Sophie Germain) \\
$3$ & $8$ & $5$ & $5/8$  & Proven (cascade moduli,~\cite{WakilEisenstein}) \\
$4$ & $16$& $12$& $3/4$  & \textbf{Open prediction} \\
$5$ & $32$& $27$& $27/32$& Conjectured \\
\bottomrule
\end{tabular}
\end{table}

\begin{remark}
Note that $k=1$ and $k=2$ both give $\theta = 1/2$. The $k=1$ case is
Bombieri--Vinogradov: one-level constitutional sieves recover exactly the
classical barrier. The $k=2$ case (Sophie Germain) is a non-instance of the
enhanced effect---the coupling prevents improvement beyond $1/2$---confirming
that the enhancement requires genuine multi-level independence ($k \geq 3$).
The first genuinely new level is $k=3$ at $\theta_W = 5/8$.
\end{remark}

%% ─────────────────────────────────────────────────────────────────────────────
\section{The Constitutional Forcing Conjecture}
%% ─────────────────────────────────────────────────────────────────────────────

\begin{conjecture}[Constitutional Forcing Conjecture]
\label{conj:cfc}
Let $\mathcal{S}$ be a system in which:
\begin{enumerate}[label=(\roman*)]
  \item the configuration space is exponential with growth rate $r > 0$;
  \item the structural law is intrinsic (not an external parameter);
  \item the system has exactly $k$ independent constitutional constraint levels,
    with $2^k$ total binary configurations of which $k$ are invalidated by the
    constraints.
\end{enumerate}
Then $\mathcal{S}$ exhibits the Shannon--Wakil effect with governing constant
$\theta^* = (2^k - k)/2^k$.
\end{conjecture}

\begin{remark}
Conjecture~\ref{conj:cfc} is deliberately broad. It encompasses information
theory ($k=1$: entropy/capacity), prime constellations ($k=3$: twin primes via
Eisenstein structure), and hypothetical $k=4$ systems yet to be identified.
The conjecture is falsifiable at every level of the $k$-hierarchy.
\end{remark}

%% ─────────────────────────────────────────────────────────────────────────────
\section{Shannon--Wakil Effect vs.\ Analogies}
%% ─────────────────────────────────────────────────────────────────────────────

It is important to distinguish the Shannon--Wakil effect from informal analogy.

An \emph{analogy} between two fields observes superficial similarity and draws
heuristic inspiration. The comparison is suggestive but not formally verifiable.
Analogies are valuable but limited: they can be wrong in ways that are hard to
detect, and they generate no falsifiable predictions.

The Shannon--Wakil effect is different in three ways.

\textbf{First}, all three instances are proved to satisfy the same formal
definition (Definition~\ref{def:sw}). This is not a qualitative resemblance but
a quantitative match: the same conditions (SW1)--(SW4) are verified in each case.

\textbf{Second}, the framework makes falsifiable predictions. The Sophie Germain
non-instance was predicted before the verification was run
(Theorem~\ref{thm:sg}); the $k=4$ prediction is open and testable
(Remark~\ref{rem:k4} below). A framework that predicts non-instances is
doing more than analogy.

\textbf{Third}, the formula correspondence (Section~\ref{sec:formula}) is
quantitatively precise: all three governing constants satisfy the same structural
formula $\theta^* = (\text{total} - \text{constraint})/\text{total}$, with the
quantities total and constraint having natural interpretations in each domain.

%% ─────────────────────────────────────────────────────────────────────────────
\section{Discussion and Open Problems}
%% ─────────────────────────────────────────────────────────────────────────────

\subsection{The 77-year gap}

Shannon proved the AEP and Channel Coding Theorem in 1948. The arithmetic
instance (cascade moduli, $\thetaW = 5/8$) was established in 2026. The
77-year gap is not surprising: the connection is visible only after both sides
are independently understood. The order---independent results first, structural
comparison second---guarantees that neither result was constructed to match the
other. This independence is what makes the parallel meaningful.

\subsection{The $k=4$ open prediction}
\label{rem:k4}

\begin{remark}[Open prediction]
\label{rem:k4}
The formula $\theta_4 = (2^4 - 4)/2^4 = 12/16 = 3/4$ predicts the existence
of a prime constellation with exactly $k = 4$ independent modular constraints,
exhibiting level of distribution $\theta = 3/4$ via the constitutional sieve
framework. Identifying such a constellation and verifying or disproving
$\theta_4 = 3/4$ is a concrete open problem that would either strengthen or
refute the $k$-hierarchy.
\end{remark}

\subsection{The Shannon--Wakil constant ratio}

The ratio $\thetaW / \thetaBV = (5/8)/(1/2) = 5/4$ is the gain in level of
distribution achieved by the Eisenstein structure over the classical Bombieri--Vinogradov
bound. Whether this ratio has a natural information-theoretic interpretation---as a
channel capacity gain from exploiting the algebraic structure of $\Zomega$ relative
to a generic channel---is open.

\subsection{Channel capacity as a forcing constant}

The channel capacity $C = \max_{p(x)} I(X;Y)$ is the maximum mutual information
over all input distributions. In the cascade moduli case, the analogue of mutual
information is the level of distribution: the supremum $\theta$ such that primes are
equidistributed over $3^K \leq x^\theta$. The parallel is:
\[
  C = \max_{p(x)} I(X;Y)
  \quad\longleftrightarrow\quad
  \thetaW = \sup\{\theta : \text{equidistribution holds for } 3^K \leq x^\theta\}.
\]
In both cases the governing constant is defined as the maximum that the structure
achieves. This maximisation correspondence is the most precise of the five and the
most promising avenue for a functorial formulation of the conjecture.

%% ─────────────────────────────────────────────────────────────────────────────
\section*{Acknowledgements}
%% ─────────────────────────────────────────────────────────────────────────────

We thank Claude Shannon posthumously. His 1948 paper established both the AEP and
the Channel Coding Theorem in a single work, and both turn out to be instances of
the same underlying phenomenon documented here. That a 1948 result in electrical
engineering is structurally parallel to a 2026 result in prime number theory would
presumably have pleased him.

%% ─────────────────────────────────────────────────────────────────────────────
\begin{thebibliography}{99}
%% ─────────────────────────────────────────────────────────────────────────────

%% ── Information theory ───────────────────────────────────────────────────────

\bibitem{Shannon1948}
C.~E.~Shannon,
A mathematical theory of communication,
\emph{Bell System Tech.\ J.}\ \textbf{27} (1948), 379--423, 623--656.

\bibitem{CoverThomas2006}
T.~M.~Cover and J.~A.~Thomas,
\emph{Elements of Information Theory},
2nd ed., Wiley-Interscience, Hoboken, NJ, 2006.

%% ── Analytic number theory ────────────────────────────────────────────────────

\bibitem{BV1965}
E.~Bombieri and A.~I.~Vinogradov,
On the large sieve,
\emph{Mathematika}\ \textbf{12} (1965), 201--225.

\bibitem{GPY2009}
D.~A.~Goldston, J.~Pintz, and C.~Y.~Y{\i}ld{\i}r{\i}m,
Primes in tuples~I,
\emph{Ann.\ of Math.}\ \textbf{170} (2009), 819--862.

\bibitem{Pascadi2025}
A.~Pascadi,
Level of distribution $5/8$ for smooth moduli,
arXiv:2505.00653v2, 2025.

%% ── Companion papers (Wakil's Program) ───────────────────────────────────────

\bibitem{WakilKhayyam}
K.~Wakil,
Khayyam's triangle: historical restoration and the geometry hidden in the
numbers,
ARC Institute of Knowware, preprint, 2026.

\bibitem{WakilSC}
K.~Wakil,
Spherical cow philosophy: a methodological framework for mathematical
discovery,
ARC Institute of Knowware, preprint, 2026.

\bibitem{WakilRamanujan}
K.~Wakil,
Ramanujan's dimensional forcing: a universal formula for sieve levels,
ARC Institute of Knowware, preprint, 2026.

\bibitem{WakilEisenstein}
K.~Wakil,
Ternary forcing on the hexagonal lattice: Eisenstein integers, cascade
moduli, and level of distribution~$5/8$,
ARC Institute of Knowware, preprint, 2026.

\bibitem{WakilW3}
K.~Wakil,
Condition~W3: absence of Siegel zeros for characters modulo~$3^K$ via
CM~structure of~$\Z[\omega]$,
ARC Institute of Knowware, preprint, 2026.

\end{thebibliography}

\end{document}
